\documentclass[11pt]{article}
\usepackage{amsmath,amssymb}
\usepackage[margin=1in]{geometry}
\usepackage{hyperref}

\title{Ideas Parked --- Not Yet Merged (v11)}
\author{Staging document; each entry marks its eventual destination}
\date{}

\begin{document}
\maketitle

\noindent This version reviews v10 against \texttt{PA\_ReWrite\_V15\_22.tex}
specifically, the version that actually carries the verb-inventory
rewrite and the Section~2/Section~3 restructuring. Three kinds of
change follow from that review, kept in separate sections below:
terminology in v10's own text that the rewrite has made stale (most of
it symbols, one of it a genuine naming collision with a new formal
operator); status markers on v8--v10 that said ``not merged'' and are
now wrong, since the realignment they describe is the current paper;
and gaps the rewrite itself opened by giving Collapse formal,
axiom-level standing, none of which were trackable before that
standing existed. Everything else in v10 was checked against the
current paper directly and found unchanged and still accurate; it is
carried forward below rather than merely cited, so this file stays a
complete, standalone record.

\section{Status update: the v8--v10 terminology work is merged, not
scribbled}

v8 proposed the realignment and flagged a symbol collision. v9
resolved that collision, provisionally, with $\circleddash$ for
UnCollapse. v10 revised that choice --- $\ominus$ for UnCollapse,
$\setminus$ for Remove --- and called itself ``ready to draft.'' All
three said, in their own closing status lines, that nothing had been
written into any version of the paper yet.

That is no longer true. Checked directly: $\circleddash$ appears zero
times in \texttt{PA\_ReWrite\_V15\_22.tex} --- v9's own proposed symbol
was superseded by v10's later choice, which is the one actually in the
paper. $\ominus$ now means UnCollapse throughout; $\setminus$ means
Remove throughout; $\uplus$ means Bundle throughout; Mix kept its own
name after a brief detour through ``MixMerge'' and back. The full
verb table v10 proposed is now Table~1 of the paper itself, with UnMix
added to it by name. Every item v10 marked ``Settled, this entry'' ---
Collapse's non-idempotence, the three convergent arguments for it, the
$N$-storage account, $A^{\langle N\rangle}$ notation, UnCollapse's
algebra at both boundary cases --- is now the paper's own stated
content, not a scribble describing an intended future state.

\textbf{[What this means for v8/v9/v10 below.]} Their technical content
was correct and is not repeated in full here; their status lines were
right when written and are wrong now. Treat v8, v9, and v10 as a
closed historical record of how the decision was reached, not as
pending work. Kept, unedited, at the end of this file for that reason
(Section~\ref{sec:closed_history}) --- their own words describe the
reasoning accurately; only their framing as future work is stale, and
this section's note is the correction, not a rewrite of what they
argued.

\section{Stale symbols: v10's own text, read against the current
paper}

The realignment changed what $\oplus$ and $\ominus$ \emph{mean}, not
just what they are called. Several of v10's earlier sections --- all
written before the realignment was settled --- use $\oplus$ for
Gather's old bowl-forming sense and $\ominus$ for Remove's old sense,
both now reassigned. Read literally against the current paper, these
lines say something false or confusing, not merely old-fashioned.
Corrected readings, section by section:

\begin{itemize}
\item \textbf{The $\mathtt{chladni}$ citation (Section~\ref{sec:gap1}
below).} The quoted sentence --- ``Cross-mode coexistence is
therefore $\oplus$'' --- is a direct quote from Paper~1, and the
paper's own words now read $\uplus$ at that exact sentence, confirmed
directly against the current text. The quote below is corrected to
match. Two more instances in that same section's commentary (``already
living on the $\oplus$ side,'' ``compositing unflattened layers is
$\oplus$'') meant the same bowl-forming sense and are corrected the
same way.

\item \textbf{Meta vectors, link-vector grouping.} ``Gathering, in the
$\oplus$ sense'' meant Bundle's collecting act. Now reads: gathering,
in the $\uplus$ sense.

\item \textbf{The v5 identity-handle scribble, most seriously.} ``Gather
still satisfies $A\oplus A=A$'' was true when written. Read against
the current paper, where $\oplus$ is Collapse, this sentence is now
\emph{false} --- Collapse's whole point, settled in v10's own later
section, is that $A\oplus A\ne A$. The claim itself was never about
Collapse; it was always about Bundle's idempotence. Corrected: ``Bundle
still satisfies $A\uplus A=A$.'' The same section's ``not collapsed by
$\oplus$'' and ``under $A\oplus A=A$'' need the identical fix, to
$\uplus$, for the identical reason. One more symbol in the same
section, ``do not claim that $\{A\}\ominus A$ ... inverts $\boxplus$,''
used $\ominus$ for Remove, which is what the sentence is actually
about; under the current symbol it reads $\{A\}\setminus A$. The
underlying claims were correct throughout and remain correct; only the
symbols were overtaken by a later decision these sections predate.
\end{itemize}

\section{A genuine naming collision, not just a stale symbol}
\label{sec:uncollapse_collision}

The v5 scribble's own name for its mechanism --- ``uncollapse,''
lowercase, informal --- is a direct collision with UnCollapse,
capitalized, now a formally defined operator with its own equation and
algebra (Section~\ref{sec:closed_history}, v10's own entry). These are
not the same thing, and the collision is not cosmetic:

\textbf{v5's mechanism:} a pointer, written at the moment a bowl is
read by Mix, from the resulting state back to the bowl that produced
it. Looking the bowl back up later via that pointer. This is
navigation by a kept record; it does not invert Mix, and v5 says so
explicitly. It is a workaround for Mix's own proven irreversibility,
built entirely from meta-level bookkeeping ($id$ fields, link
vectors) that sits outside the 4-tuple.

\textbf{The paper's UnCollapse:} removes one informational vector from
a Collapsed atom's own internal list, $N\to N-1$, by direct
computation on the atom's own structure --- no pointer, no bowl, no
meta field, nothing external. It answers a different question (what
contribution is taken back from a count-preserving combination) using
different machinery (Axiom~7's own vectors) for a different operator
(Collapse, not Mix).

\textbf{Why this matters going forward.} v5's mechanism is still a
live, useful idea --- a real answer to ``can anything be done about
Mix's irreversibility,'' and nothing above withdraws it. But it must
never again be called ``uncollapse.'' Proposed name, used from here on:
the \textbf{Mix-provenance link}. $id$ stays the field name (unaffected
--- it was never a verb, so it does not collide); only the informal
verb ``uncollapse'' retires, in favor of ``recover via Mix-provenance''
or similar. The tiny worked picture from v5 (bowl $\{$red,
green$\}$, Mix reads yellow, provenance link from yellow's $id$ back to
the bowl's) is unaffected in substance and is restated below under the
corrected name.

\section{Three gaps the rewrite opened, not yet tracked anywhere}

None of these existed as trackable items before Collapse had formal,
axiom-level standing. All three are visible directly in the current
paper's own text --- two as explicit, stated omissions; one as a
silence worth naming rather than a stated one.

\paragraph{Collapse across differing sorts (the paper already says
this is open).} Quoted directly: ``Axiom~7 currently requires every
vector in one atom to share the atom's single $\tau$ ... Admitting a
per-vector $\tau_i$ would be the natural extension ... but it is not
adopted here and is left open.'' This is Table~2's own open$^\P$
entries for every cross-sort pair. Worth tracking here specifically
because it is the same shape of question as Gap~1
(Section~\ref{sec:gap1}) for a different operator, and the two should
not be assumed to share an answer: Mix blends, so refusing across
directions protects against manufacturing a direction that was never
there; Collapse preserves every contribution without blending
anything, so the case for refusing it across sorts is not the same
case, and has not been separately made.

\paragraph{Collapse as $T$'s actor: not claimed, not ruled out.}
Quoted directly: ``every worked $T$ path in this paper takes one
[teaching atom], and a Collapsed-atom-as-actor example is not claimed
here.'' $T$'s machinery for accepting a richer actor already exists in
principle --- $V_{\text{derived}}$, or beyond three dimensions,
$N_{\text{dim}}-2$ of the Collapsed atom's own orthonormal arrows
directly --- but no example walks a Collapsed atom, and the cast that
would need naming at the call site is not worked out anywhere.

\paragraph{Assembly and Collapse: no stated relationship at all.}
Checked directly: Axiom~6 (Assembly) does not mention Collapse once,
in either direction. Whether a Collapsed atom may stand as one of
Assembly's two related units, whether $U_{AB}$ means anything when $A$
or $B$ carries $N>1$ internal vectors, is not addressed, not refused,
not flagged as open --- simply silent. Of the three gaps here, this is
the one with no sentence in the paper pointing at it; recorded so it
is not mistaken later for a settled ``no.''

\section{A new route worth recording, not a solution}

The ``reducing a path by merging close actors'' open item
(Section~\ref{sec:pathreduction} below, unchanged and still open in the
current paper) asks, word for word, whether a representative actor
should be ``a $\boxplus$-blend of the close actors, an average
rotation, or something else.'' Collapse did not exist as a named
option when that sentence was written. It is a structurally plausible
candidate for ``something else,'' worth recording precisely because of
how it would differ from the other two: Collapse the close actors
losslessly first --- every contribution kept, nothing blended, nothing
decided yet about direction --- and only then define a second, reading
operation (the way Mix already reads a Bundle) that would extract a
single representative direction from the Collapsed result.

This is not a solution. The reading step that would turn a Collapsed
atom into one representative direction does not exist, is not sketched
here, and may turn out to face the same direction-blending objection
that keeps Gap~1 open --- extracting one direction from several is
exactly the operation Route~1's chladni evidence counsels refusing.
What Collapse changes is only that the open item's own ``or something
else'' now names something specific, where before it named nothing.
The same observation extends, with the same limit, to the v6
item (Section~\ref{sec:v6} below): for same-sort regions it offers an
alternative to the already-working Hill bijection, not a need; for
cross-sort regions it offers nothing new, since Collapse cannot yet
cross sorts either (the first gap above) --- the hard half of v6 stays
exactly as open as Gap~1 itself.

\section{Destined for Pattern Arithmetic (still formally open ---
unchanged, confirmed against the current paper)}
\label{sec:gap1}

\paragraph{Mix across differing directions.} Checked directly: the
paragraph's own text in \texttt{PA\_ReWrite\_V15\_22.tex} still reads
``not settled here,'' still lists the same three unchosen routes,
still carries the same withdrawn provisional recipe. Nothing here
changed it, on explicit instruction to leave gap work alone until
terminology settled; recorded accurate and current, not merely
inherited.

\textbf{[Found: the paper already answers this question once, for one
sort, and points toward a general answer.]} The $\mathtt{chladni}$
inclusion states, word for word (corrected symbol, per
Section~\ref{sec:uncollapse_collision}'s sibling correction above):
``Across modes the map is not available, and not for want of effort.
Mode space is not closed under addition... Cross-mode coexistence is
therefore $\uplus$ --- a bowl of modes --- and the resulting figure is
a reading of that bowl which is not $\boxplus$. This is the open
question of [the differing-directions open item] met in a concrete
domain, and it suggests the question may be closed negatively for this
sort rather than left waiting.'' This is a citation, not a new
argument --- the paper's own words, in a section that was not
previously read as bearing on Gap~1 directly.

\textbf{[Why this converges with everything else worked through in
this project, not just chladni.]} Once Mix is read as ``resolve
several reports on one shared question into a verdict'' (argued in
v2/v3, carried forward unchanged), refusing across-direction
combination stops looking like a gap and starts looking like the
correct behavior of an operator that was never asked that question in
the first place. And every practical workaround that actually worked
in this project was already living on the $\uplus$ side of precisely
this line, independently: bowls holding separate weather channels,
per-channel similarity scoring, the candlestick summary for
reference-atom building, Assembly's own camp-to-system Hill bijection.
None of them ever needed Mix to cross a direction. All four arrived,
separately, at ``keep it in a bowl, read it there'' --- which the
chladni section had already stated as a formal result.

\textbf{[Status: a leading candidate, not yet a decision.]} This
supports Route 1 (refuse outright, from the three already on file)
over Route 2 (rotate first) or Route 3 (two-stage operator) as the
likely resolution, with a real citation behind it rather than
preference. Writing this into Paper 1 would mean editing the
``Open: Mix across differing directions'' paragraph to state the
refusal explicitly, citing the chladni section as precedent. Not done
here, on request --- parked, ready, pending a decision to merge.

\textbf{[Carried forward unchanged: the reading argued in v2/v3.]} Mix
is not ``multiple vectors becoming one vector.'' It takes several
reports already agreed to answer one shared question and returns a
verdict about that question. The shared $\hat n$ is the precondition
that makes ``these answer the same thing'' true, not a boundary
condition on a more general operator. Supporting evidence: same-
direction opposite-phase cancellation to black (principled
contradiction-handling, not a fabricated compromise); unmixing's
proven ill-posedness, with metamerism as independent real-world
precedent; the paper's own opening line, ``two notes yield a chord,
not a larger integer... the word `add' names a different rule.''

\textbf{[Carried forward unchanged: the ``we mix all the time''
challenge, resolved.]} Sugar in tea is one dial, never needed
differing directions. Red plus green additively sits ``at the same
place'' in the paper's own words --- already inside Mix's settled
domain. Paint-layer flattening is confirmed, across every source
checked, as destructive and irreversible, matching unmixing exactly;
compositing unflattened layers is $\uplus$, not $\boxplus$, which is
why ``unmerge them later'' was always possible. Representing a whole
painting as one vector remains Gap 1 itself, not a counterexample to
it.

\textbf{[Carried forward unchanged: the two-part split, now read as
one part likely answered.]} Embedding validity --- does geometric
closeness on a shared sphere track real semantic agreement at all ---
was the harder, prior question, true for hue by deliberate
construction, unstated for an arbitrary sphere of concepts. Given the
chladni precedent above, the likely answer for sorts where this was
never stated is that no valid embedding needs constructing, because
cross-direction combination is refused before the embedding question
would even matter.

\textbf{[Carried forward unchanged: still confirmed not to help.]}
Magnitude self-bounding near opposition does not license the
direction that comes with it. The Hill form and Assembly's Hill
bijection both fix magnitude only, verified twice, never direction.

\section{A correction, recorded so it is not repeated}

\paragraph{$\tau_{\text{weather}}$ names two different problems, and
only one of them is Gap~1.} Raised once as a proposed admissible sort
for combining different weather concepts (rain, drought) across
directions --- that question is genuinely Gap~1, and the discussion
above bears on it. Raised a second time, later, meaning something
narrower: given several already closely-aligned vectors within one
bounded region, how to arrive at a single representative atom, using
the Hill form. That second question does not invoke $\boxplus$ across
differing directions at all --- it is Assembly's already-verified
system-promotion mechanism (16 camps to 13 systems under the Hill
bijection), recognized in a new setting, not a Paper~1 admissibility
question. The two were briefly conflated in this project; recorded
here so the same conflation is not repeated. Assembly's version needs
nothing further; Gap~1's version is addressed in
Section~\ref{sec:gap1} above. Unaffected by the rewrite; checked, no
terminology in this section needed any update.

\section{Destined for Pattern Arithmetic (advances an existing open
item, unchanged)}

\paragraph{Location as a sort --- a third candidate.} The existing
paragraph names two candidates and adopts neither. Confirmed present
and unchanged in the current paper: \texttt{PA\_ReWrite\_V15\_22.tex}'s
own ``Open: location as a sort'' paragraph still asks the identical
question in identical terms. A third was tested here: dedicate
separate, disjoint dimensions to content, to time, and to location,
rather than letting any of them share an axis.

\textbf{[Verified.]} Two shifts acting through disjoint content
dimensions commute exactly (checked numerically, not approximately:
time-then-location and location-then-time landed on the identical
point to full floating-point precision). Two shifts forced to share a
single content dimension do not commute, confirming the disjointness
is what matters, not rotation itself.

\textbf{[Verified.]} Content stays recoverable after both shifts:
projecting back onto the content dimensions alone reproduced the
original direction closely, while the time-axis and location-axis
each carried a clean, separately-readable amount of shift, with
nothing blended between them.

\textbf{[Open, unchanged.]} The encoding itself --- what specific
rotation a given time value or location corresponds to --- is not
stated anywhere. The geometry now supports it cleanly; the mapping
does not exist yet.

This is real, checkable content, not yet written as an addition to
the existing paragraph.

\section{Destined for Pattern Arithmetic (open item stays open,
unchanged, now with a newly available route noted above)}
\label{sec:pathreduction}

\paragraph{Reducing a path by merging close actors.} The existing
paragraph asks for three things: a stated distance, a merge rule, and
a bound on how far a reduced path's endpoint may drift from the
unreduced one. Discussion here supplied a distance (see below) and
partially addressed the merge rule, but only for magnitude, not for
direction --- and direction is the harder half this item is actually
about. Confirmed present, unchanged, in the current paper; see
Section~5 above for what its existence newly suggests, not yet a
resolution.

\textbf{[Verified.]} Switching the strength calculation from the
companion paper's clip to Inference's already-established Hill form,
$r'=|z|/(K+|z|)$, changes nothing about the resulting direction
$\chi'=\arg(z)$. Checked directly on two constructed sequences with
identical content and opposite order: both gave the same angle under
either formula. The Hill substitution fixes evidence-count blindness
in strength; it does not touch the order-blindness this open item is
actually asking about.

\textbf{[New: a candidate grounding for the stated distance,
item~(1) of the three things required.]} How many actors can be
mutually distinguishable --- separated by at least some angle
$\theta_{\min}$ --- on $\mathcal{S}$ at all is a real, named problem
(spherical codes, the Tammes problem), not an arbitrary choice to
make. Exact optimal counts are known only for small $N$; a standard
approximate bound, treating each point as owning a non-overlapping
cap of half-angle $\theta_{\min}/2$, gives
\[
N_{\max} \;\approx\; \frac{4\pi}{2\pi\bigl(1-\cos(\theta_{\min}/2)\bigr)}
\;=\; \frac{2}{1-\cos(\theta_{\min}/2)}.
\]
Computed for illustration:
\begin{center}
\begin{tabular}{rr}
$\theta_{\min}$ & approx.\ max.\ distinct points \\
\hline
$1°$   & $52{,}525$ \\
$5°$   & $2{,}101$ \\
$10°$  & $526$ \\
$30°$  & $59$ \\
$60°$  & $15$ \\
$90°$  & $7$ \\
$120°$ & $4$ \\
\end{tabular}
\end{center}
This is an approximation, not an exact extremal result, but it is a
real, checkable mathematical object rather than a guess: naming a
$\theta_{\min}$ for ``close enough to merge'' directly bounds how many
genuinely distinct actors can exist near each other at all, which
bears on item~(3), the drift bound, as well as item~(1).

\textbf{[Open, unchanged.]} No merge rule for direction exists. No
drift bound exists for either magnitude or direction. This item
remains exactly as open as it was before this discussion, now with
the magnitude half more precisely separated from the direction half,
a real candidate supplied for the distance half specifically, and one
newly nameable candidate mechanism for the merge rule itself
(Section~5 above), not yet worked out.

\section{Destined for Inference (new mechanism, unchanged)}

\paragraph{Candlestick summary: a richer merge rule for reference-atom
building.} Not a resolution of the Paper 1 item above --- a different
question. Paper 1's open item is about \emph{walking}: can a shortened
path still be walked to roughly the same endpoint. This is about
\emph{reading}: how to summarize a group of atoms for comparison
against a reference atom, which is Inference's own established
business (its reference-atom section), not $T$'s. Destined for a
different paper; unaffected by anything in this review.

\textbf{[Verified.]} Open/close/high/low, computed the way a trading
candle is, distinguishes exactly the case flat Mix could not. Two
sequences, identical content, opposite order --- naive Mix gave the
identical angle ($47.5°$) for both. The candle's open and close swapped
between them ($20°\!\to\!160°$ versus $160°\!\to\!20°$), correctly
reading one as ending badly and the other as ending well. High and low
stayed identical for both, correctly capturing that only the
\emph{range} was shared, not the \emph{direction of travel}.

\textbf{[Settled, matches existing machinery.]} The scope rule ---
every atom belongs to exactly one time unit, no overlap between
buckets --- is not new; it is precisely the discrete-time-bucket
metadata already defined on a walk's own steps in the companion paper.

\textbf{[Open.]} This gives Inference a genuinely better reference-atom
merge rule than flat Mix. It is not proposed as a fix for Paper 1's
walking-specific open item, which remains untouched above. Whether
open/close/high/low itself needs a coherence check analogous to
$\bar R_c$ has not been considered.

\section{Destined for Inference (naming caution, unchanged)}

\paragraph{``Inference Mandelbrot'' is not Assembly's Mandelbrot, and
should not share its name.} Assembly's multi-zoom promotes an
unordered thing (a settled camp) and is fully handled by the Hill
bijection already verified there. Whatever would promote an
\emph{ordered} thing (a walked path, a movie, a time series) into one
higher-level atom is a different, harder problem --- it is the
direction half of the path-reduction item above, restated from a
different angle. \textbf{[Naming flag only.]} If this is ever built,
it wants its own name; reusing ``Mandelbrot'' for it would be the same
kind of collision already caught and corrected for ``collapse,''
``singularity,'' and the quantum swap test --- and now, per
Section~\ref{sec:uncollapse_collision} above, for ``uncollapse'' as
well.

\section{Destined across three papers (new, split by piece; symbols
corrected, substance unchanged)}

\paragraph{Meta vectors: a second kind of vector alongside information
vectors.} Proposed: an information atom may carry two distinct kinds
of vector. Information vectors are Axiom 7's bundle $v_i$'s, capped at
$r_i\le1$, potentially several sorts $\tau$ across a whole record.
Meta vectors describe the atom rather than being part of its content
--- three named so far: $V_{\text{derived}}$, link vectors (pointers
between atoms, one- or two-directional), and a timestamp vector.
Explicitly left open whether more meta-vector types are needed. Meta
vectors are not capped at 1, and grouping several atoms' meta vectors
together on like terms is proposed as a distinct mechanism for
building higher structures. This is one idea that lands in three
different places once traced to where each piece actually belongs.

\textbf{[Destined for Pattern Arithmetic.]} The taxonomy itself ---
information vectors versus meta vectors as a formal, named distinction
sitting on top of Axiom 7 --- is structural, the same kind of addition
Axiom 7 itself was. $V_{\text{derived}}$ already lives there. A
timestamp vector, if adopted, would be a third, genuinely different
treatment of time from the two already on file above: distinct from
the discrete-bucket metadata already defined on a walk's own step, and
distinct from the dedicated-axis encoding in the location-as-a-sort
entry. Not yet clear whether more than one of the three should
coexist, or which.

\textbf{[Destined for Assembly.]} Link vectors specifically. This is
the mechanism the shelf's own neighbour-graph note was already naming
without a concrete form: ``a data structure on seats after rest, not
$T$.'' Grouping link vectors on like terms is collecting edges into a
graph, not merging them --- gathering, in the $\uplus$ sense --- so
it carries none of the open questions below it.

\textbf{[Destined for Inference.]} The grouping mechanism itself:
querying several atoms' meta vectors, confirming they sit close
together, and promoting one shared meta-vector up to the group as a
whole. This is the space-time gating idea from earlier discussion,
generalized into a standing principle. Timestamp and location get
\emph{used} here; they are only \emph{defined} in Pattern Arithmetic
above.

\textbf{[Verified, carried forward as a caveat rather than a new
finding.]} Combining several atoms' $V_{\text{derived}}$'s stays safe
only under the constraint the original already carries: auxiliary,
never promoted to a trusted, walkable atom on its own. Checked
directly that neither Inference's Hill form nor Assembly's Hill
bijection touches this constraint --- both fix magnitude only, never
direction --- so grouping $V_{\text{derived}}$'s inherits that same
limit rather than escaping it.

\textbf{[Renamed here, per Section~\ref{sec:uncollapse_collision}
above.]} Link vectors were separately discussed as one possible
mechanism for a limited, honest form of Mix-provenance recovery: not
inverting Mix (still impossible, unaffected by anything here), but
keeping an explicit, external pointer from a Mixed result back to the
bowl that produced it, created at Mix-reading time, before the bowl
itself might otherwise be discarded. This is navigation via a kept
record, not reconstruction from the Mixed result --- the distinction
the white-light/prism discussion turned on, confirmed there to
matter. No longer called ``uncollapse''; see
Section~\ref{sec:uncollapse_collision} for why.

\section{Destined across three papers (scribble, v5; renamed and
symbol-corrected, substance unchanged)}

\paragraph{Identity handle as a meta field, for Mix-provenance
recovery.}
An information atom may carry one more meta quantity: an opaque
identity $id$, assigned before any Bundle or Mix, and kept on the
atom afterwards. The scribble that prompted this note called it a
``meta vector'' with $\tau_{id}$ equal to that identity. Fine-tuning
below refuses $\tau_{id}$ as a sort and keeps $id$ as a \emph{field},
in the same shelf as timestamp, mint-location, and link pointers
already named in the meta-vector entry.

\textbf{[What it is for.]} Mix remains irreversible (Paper 1 v1.1).
Mix-provenance recovery here does not reconstruct ingredients from a
yellow. It means: at Mix-reading time, write a pointer from the
resulting atom (or from the register occupant after a walk) back to
the bowl that produced it, and keep that bowl under the same $id$ (or
a pair $result\_id \leftrightarrow bowl\_id$). Later, given the
result, look the bowl up. Navigation by a kept record, not inversion.
This is the same distinction already recorded under link vectors and
the white-light / prism discussion.

\textbf{[Viability --- what works.]}
An opaque handle is the right shape. It does not need to live on
$\mathcal{S}$, does not need $r\le 1$, and does not need a direction.
A string or integer key is enough. Creating the pointer \emph{before}
the bowl is discarded is the whole trick; after discard there is
nothing to recover. Links already named that graph; $id$ is the
node name those edges need. Assembly can own the edge table
(neighbour / parent-bowl). Paper 1 only needs the taxonomy: $id$ is
meta, not a unit of truth. Inference may use $id$ to re-open a bowl
and run a second reading.

\textbf{[Viability --- what to refuse now.]}
Do not mint $\tau_{id}$. $\tau$ is the legality tag of an information
vector. An identity is not a pattern and must not appear in the type
table. Do not encode $id$ as an angle on the sphere; colliding two
ids geometrically would be a false Mix. Do not claim that
$\{A\}\setminus A$ or $T^{-1}$ plus $id$ inverts $\boxplus$. Bundle
still satisfies $A\uplus A=A$: two tokens of the same pin share one
address and, if they are the same \emph{instance}, one $id$; two
distinct instances of the same pin (two physical chairs) need two
ids and are not Bundled by $\uplus$ unless the lexicon says they
are the same state. That instance-versus-kind cut is the fine-tuning
this scribble does not finish.

\textbf{[Tiny picture.]}
Bowl $B=\{\mathrm{red},\mathrm{green}\}$, $id_B=\#17$.
Mix reads yellow, $id_Y=\#18$, meta link $\#18\to\#17$.
Looking up $\#18$'s provenance returns $B$, still red and green,
still two atoms. The yellow is not pulled apart. The prism was kept
on the shelf.

\textbf{[Status.]} Scribble. Destined: taxonomy to Paper 1 (meta
shelf, next to $V_{\mathrm{derived}}$); the pointer table to
Assembly; lookup-and-reread to Inference. Not merged. Fine-tune
later: who assigns ids, instance vs kind under $A\uplus A=A$, and
whether one atom may keep a stack of parent-bowl ids after several
Mixes.

\section{A theoretical problem, split the same way an earlier section
already split one --- v6, unchanged except for the new route noted
above}
\label{sec:v6}

\paragraph{Can every atom in a unit space, across all labels and
types, be merged to one representative atom?} Posed precisely: a
bounded region holds many already-settled atoms, of possibly
\emph{differing} sorts $\tau$, not just differing content within one
sort. Can the whole collection become one representative for that
region? Framed, as requested, with both a theoretical and a practical
answer in view, the way the sphere-packing question above was.

\textbf{[Solved --- but only for the restricted case, which is an
earlier correction's general form.]} When every atom in the
region shares one sort (or, more precisely, one already-admitted
$\varphi_\tau$ direction-family), this is exactly
Assembly's already-verified system-promotion mechanism: many close
atoms, one representative, via the Hill bijection rather than the
naive clip, because Hill preserves evidence-count where the clip
saturates and loses it. Sixteen camps correctly became thirteen
systems on this basis, not eleven. ``A lot of random atoms'' changes
nothing here --- randomness in position or content is fine, provided
the sort is shared. This is the $\tau_{\text{weather}}$
correction, restated as the general pattern it always was: random
instances of one $\tau$, closely aligned, merge cleanly. Section~5
above notes Collapse as an alternative route for this same-sort case
specifically --- not a need, since Hill already works here, but a
newly available option.

\textbf{[Still open --- this is Gap~1, not a second result.]} ``Across
all labels and types,'' taken literally, names the harder case: atoms
of genuinely \emph{different} sorts in the same region. This is not
solved by Hill's equation or the bijection, and --- confirmed by
Section~4 above --- not solved by Collapse either, since Collapse
cannot yet cross sorts for exactly the same reason Mix cannot.
Confirmed twice already in this project: both Inference's Hill form
and Assembly's Hill bijection fix magnitude saturation only, and
neither touches direction or sort. Merging across genuinely different
sorts into one atom is precisely Gap~1 (Section~\ref{sec:gap1}), for
which a leading candidate (refuse outright, per the $\mathtt{chladni}$
precedent) exists but has not been written into Paper~1.

\textbf{[Theoretical versus practical, for the cross-type half
specifically.]} The sphere-packing question above had a well-defined
continuous quantity (point count) that a practical approximation could
meaningfully estimate. This question's open half does not: there is no
stated $\varphi_\tau$ for combining across sorts to approximate in the
first place. A random or approximate merge rule cannot stand in for a
merge rule that has not been defined --- this is an admissibility gap,
not a precision gap, and the two should not be read as the same kind
of ``open'' question.

\textbf{[So, precisely, correcting the framing in which this was
posed.]} ``We solved it using bijection and Hill's equation'' is true
for same-sort regions, and is a real, verified, reusable result ---
but it is not yet true for atoms spanning genuinely different types,
which remains exactly Gap~1, unresolved by anything checked so far.
The two cases should be kept as separate claims, the same way an
earlier section already separated them once for $\tau_{\text{weather}}$
specifically; this entry generalizes that same correction rather than
adding a new one.

\section{Destined for Pattern Arithmetic (scribble, v7; unchanged, not
yet merged)}

\paragraph{``Refuse outright'' is a claim about $\boxplus$, not a
permanent bar on ever combining unlike sorts by any mechanism.}
Raised as a caution on Gap~1's leading candidate: likely that only
like-typed atoms will ever practically need merging, but unwise to
state Route~1 in a way that forecloses future, explicitly-justified
mechanisms for binding unlike sorts, should a practical need for one
arise. Checked against the current paper: still not written into the
``Open: Mix across differing directions'' paragraph; still exactly as
scoped a caution as when drafted.

\textbf{[The distinction that matters.]} Route~1 was always a claim
about $\boxplus$ specifically --- that this particular operator, as
defined, was never asked the cross-direction question and correctly
declines to answer it. It was never a claim that no mechanism, under
any name, could ever combine unlike sorts. Conflating the two would
overclaim something the evidence never supported.

\textbf{[Why this is not a new principle, but a confirmation of one
already load-bearing in the paper.]} Axiom~0 is already built this
way: ``a sort absent from this table is undecided, not admitted;
treat it as inadmissible until a $\varphi_\tau$ is stated.''
Inadmissible \emph{now}, for lack of a stated proof --- not
inadmissible forever, by declaration. $\mathtt{chladni}$ is the
existing demonstration: admitted later, on stated grounds, not
pre-emptively barred. Even $\mathtt{chladni}$'s own hedge on
cross-mode combination, the strongest evidence behind Route~1, uses
``may be closed negatively... rather than left waiting'' --- a
leaning, not a verdict.

\textbf{[How Route~1 should actually be worded, if merged.]} Not:
``$\boxplus$ refuses states at differing directions, full stop.''
Better: ``$\boxplus$, as currently defined, refuses states at
differing directions; a cross-sort reading remains possible in
principle, as its own stated mechanism, admitted on its own terms the
way any new capability is --- not as a modification of $\boxplus$
itself.'' This keeps the refusal honest about what has actually been
shown without overclaiming permanence the evidence does not support.

\textbf{[Status.]} Scribble. A scoping note on Gap~1's existing entry,
not a new result and not yet written into that section's own text.
``Probably rare'' (practical likelihood) and ``permanently barred''
(structural impossibility) are different claims; only the first has
support, and Route~1's eventual wording should reflect only that one.

\section{Closed history: v8, v9, v10, kept verbatim}
\label{sec:closed_history}

The three sections below are reproduced unedited from v10. Their
technical content is correct and is now the paper's own content, not
a proposal; their own status lines (``scribble,'' ``not merged,''
``awaiting a decision'') describe their state at the time of writing,
not the state today. See Section~1 above for the correction. Kept
here, word for word, as the historical record of how the current
verb inventory was actually reached --- useful for exactly that
reason, and for no other: nothing below should be read as open work.

\subsection*{v8 --- proposed terminology realignment, with cost
assessment (historical)}

\paragraph{Proposed: Collapse, UnCollapse, Bundle, MixMerge, UnMix.}
A full realignment of five terms, proposed together:
\begin{itemize}
\item \textbf{Collapse} ($\oplus$) --- one atom, one origin, several
      internal vectors.
\item \textbf{UnCollapse} ($\ominus$) --- subtract one atom's vectors
      from another; not previously discussed in depth.
\item \textbf{Bundle} --- a bowl/set, many atoms, kept separate,
      separate origins.
\item \textbf{MixMerge} ($\boxplus$) --- one new atom, fused.
\item \textbf{UnMix} --- not possible.
\end{itemize}
This reassigns Collapse and Bundle relative to the reading settled
two exchanges earlier in this project (Collapse $=\oplus=$ the bowl,
confirmed three times directly from source text and from the bowl's
own definition as a separate finite set). The proposal here reverses
that assignment. Recorded as a deliberate choice to evaluate, not
adopted, and not written into Paper~1.

\textbf{[Blocking: a direct symbol collision.]} $\ominus$ is not
available for UnCollapse. It already carries its own formal axiom:
``Axiom 4: Gather and remove ($\oplus$, $\ominus$)... $\ominus$
undoes a gather. It is not anti-phase interference.'' Two
incompatible operators cannot share one symbol in the same document.
A different symbol is needed before UnCollapse can be merged under
any sequencing; this is not a matter of care in editing, the symbol
itself is already spoken for.

\textbf{[Instance counts, checked directly against the source.]}
\begin{center}
\begin{tabular}{llcl}
Term & Current meaning & Count & Title usage \\
\hline
Collapse & $\oplus$'s bowl-forming act & 14 & Yes --- a subsection
title \\
Mix / $\boxplus$ & fused, single output & 77 / 50 & Yes --- Axiom~5's
own title \\
Bundle & one atom, $N$ internal vectors & 10 & No \\
Bowl & a set of separate atoms & 80 & No, used throughout \\
\end{tabular}
\end{center}

\textbf{[Cost, per piece.]}
\begin{itemize}
\item \emph{Collapse as one atom:} high cost, and a reversal rather
      than a rename. The ``Process principle: Collapse, path,
      result'' subsection's own stated reasoning --- ``collapse is
      the act that discards the order and count of that sequence...
      $R=$ the bowl'' --- would need rewriting, not relabeling: its
      current logic explains why collapse yields a set, and that
      explanation would need to become its own opposite.
\item \emph{Bundle as bowl:} high cost. Collides with Bundle's own
      existing, different meaning (Axiom~7's $N$-vectors-in-one-atom).
      Bowl already carries 80 instances of its own, separate meaning;
      this is a reconciliation of two already-distinct, already-used
      terms, not an added synonym.
\item \emph{MixMerge, rename only, same underlying behavior:}
      moderate cost, mechanical. 77 instances is substantial, but
      nothing about $\boxplus$'s actual behavior changes --- risk is
      missed instances or inconsistent phrasing, not logical
      breakage.
\item \emph{UnMix, named but still impossible:} low cost. The concept
      is already stated (``division as unmixing is shown ill-posed'')
      --- this mostly formalizes an existing result with a name and a
      symbol, once a free symbol is chosen.
\end{itemize}

\textbf{[Overall, honest verdict.]} The $\ominus$ collision is a hard
blocker, independent of sequencing, and needs resolving before
anything else proceeds. Past that, Collapse-reversal and
Bundle-reassignment are the expensive, high-risk pair, each touching
foundational, load-bearing text and interacting with one another.
MixMerge and UnMix are comparatively cheap and could move first,
independently, as a lower-risk starting point if this is pursued.

\subsection*{v9 --- symbol resolution for Bundle and UnCollapse, plus
when ids are actually needed (historical)}

\paragraph{Bundle gets a symbol and becomes a verb.}
\[
\text{Bundle}:\quad A_1 \uplus A_2 \uplus \cdots \uplus A_n
\;\longrightarrow\; B \qquad (\text{many atoms, kept separate}).
\]
Checked clean against the source: $\uplus$ (multiset union) does not
appear anywhere in Paper~1. The choice is principled, not arbitrary
--- $\uplus$'s standard mathematical meaning, combine while
preserving separate membership, is exactly Bundle's job. Worth being
precise about what this move actually is: with $\oplus$ reassigned to
Collapse under the v8 scheme, Bundle-as-verb takes over the role
Paper~1's own $\oplus$ currently performs (the bowl-forming act,
nothing merged, nothing lost). Not a new mechanism --- a new name and
symbol for an existing one, freed up because the old symbol moved.

\paragraph{UnCollapse gets a symbol.} $\circleddash$, also checked
clean against the source. Visually related to $\ominus$, so the
family resemblance between Collapse/UnCollapse reads correctly, while
remaining a genuinely distinct symbol --- no collision with Axiom~4's
own $\ominus$. \emph{[Superseded by v10 below: the symbol actually
adopted is $\ominus$ itself, not $\circleddash$, once Remove moved to
$\setminus$.]}

\textbf{[Precisely when ids are needed for UnCollapse, and when they
are not.]} Collapse, under this scheme, is lossless by construction:
every $v_i$ stays individually intact inside the resulting atom.
Splitting it back apart is \emph{already} structurally possible
without ids whenever every incoming vector has a distinct
$(r_i,\hat n_i,\chi_i,\tau_i)$ --- the bundle's own listing already
carries everything needed.

Ids earn their place in exactly one case: two \emph{distinct} source
atoms that happen to share identical values. Without an id, the two
resulting entries inside the bundle are indistinguishable from each
other --- UnCollapse would return the right values, with no way to
confirm whether this was genuinely two separate instances that
happened to match, or something else that produced the same numbers.
This is the same instance-versus-kind distinction already on file
from the v5 scribble: ``two distinct instances of the same pin...
need two ids.''

So, stated precisely rather than as a blanket claim: UnCollapse is
value-exact with or without ids. What ids add is identity-exact
recovery --- confirmation that what comes back out is traceably the
same atom that went in, not merely one with matching numbers. For the
ordinary, all-distinct case, this guarantee is free. For the
degenerate, duplicate-value case, it does not exist without an id.
That is the honest scope of ``may become a possibility with ids'' ---
a fix for one specific, real ambiguity, not a general precondition for
UnCollapse to work at all. \emph{[This precise scoping is the one
piece of v9 directly visible in the current paper, cited by name in
UnCollapse's own section.]}

\subsection*{v10 --- the full verb inventory settled (historical;
this is the version that was drafted)}

\paragraph{Final symbol assignment, reasoned rather than aesthetic.}
Proposed: $\ominus$ moves from Remove to UnCollapse; Remove takes
$\setminus$ instead. Not merely a preference --- $\oplus$ and
$\ominus$ are a matched shape-family by standing convention (same
circle, opposite sign), and that pairing belonged to Gather/Remove
only while $\oplus$ meant Gather. Now that $\oplus$ means Collapse,
the natural partner follows it: $\oplus$/$\ominus$ as Collapse/
UnCollapse restores the convention rather than breaking it. $\uplus$
(Bundle) never had a natural circled-minus partner; $\setminus$, set
difference, fits it better and matches Bundle's own description as
practically a set operation. Checked clean against the source: zero
instances of \texttt{\textbackslash setminus} anywhere in Paper~1.

\textbf{Final table:}
\begin{center}
\begin{tabular}{llll}
Verb & Symbol & Produces & Status \\
\hline
Bundle & $\uplus$ & bowl, atoms kept separate & Settled \\
Remove & $\setminus$ & undoes Bundle & Settled (was $\ominus$) \\
Collapse & $\oplus$ & one atom, count-preserving & Settled, this entry \\
UnCollapse & $\ominus$ & removes one contribution & Settled, this entry \\
MixMerge & $\boxplus$ & one atom, fused & Settled (rename only) \\
UnMix & $\boxminus$ & --- & Proven impossible, Section~7.2 \\
\end{tabular}
\end{center}

\paragraph{Bundle's idempotence: confirmed, already proven, nothing
new.} $A\uplus A=A$. Not a new claim --- stated, word for word, under
the retiring name: ``$\oplus$ collects: identical tokens do not
double ($A\oplus A=A$)... $\{A\}\cup\{A\}=\{A\}$.'' Ordinary set
theory already backs this; a set never had a duplicate-counting
mechanism to begin with, so $\{A_1,A_1\}$ and $\{A_1\}$ were never
different objects. Confirmed by direct derivation from the ``asleep
on the key'' case: many accidental repeats of one atom, Bundled,
still return the same single bowl entry.

\paragraph{Collapse's idempotence: resolved as genuinely different
from Bundle's, after real back-and-forth.} $A\ \text{Collapse}\ A \ne
A$ in general. Reached by three convergent arguments, not one:

\begin{enumerate}
\item \emph{Information-loss.} Collapse's own stated purpose is that
no information is lost. If $A\ \text{Collapse}\ A=A$, the second
contribution is discarded exactly the way Bundle already discards
duplicates --- which would make Collapse behave identically to
Bundle in this one case, undercutting the reason the two needed to be
separate verbs.
\item \emph{Algebraic contradiction.} If $A\ \text{Collapse}\ A=A$,
then UnCollapsing one $A$ back out must return the other $A$, giving
$A\ \text{UnCollapse}\ A=A$ --- which directly contradicts
$A\ \text{UnCollapse}\ A=0$ (the atom-level analogue of
$\{A\}\ominus A=\emptyset$, independently solid). Both cannot hold
unless $A=0$.
\item \emph{The evidence-count parallel.} This is the same shape of
failure that motivated fixing Mix's cap with the Hill form: if exact
duplicates are silently discarded, two independent contributions that
happen to agree exactly become indistinguishable from one
contribution stated once --- precisely the count-blindness that was
expensive to fix once already, now reappearing in a different verb.
\end{enumerate}

\textbf{[One self-correction recorded honestly.]} An early argument
for this conclusion leaned on $T_w\circ T_w\neq T_w$ (repeated visits
in a path are genuine, separate steps) as a direct analogy. That
analogy was flagged and partly withdrawn: a path's repeated visits
are sequential, temporally separate events, while Collapse's multiple
contributions are simultaneous and co-present --- not obviously the
same kind of multiplicity. The three arguments above do not depend on
that analogy and are treated as the actual basis for the conclusion.

\paragraph{Where $N$ lives, and why nothing new needs storing.} $N$
is not a new field. Axiom~7 already defines it as the length of the
bundle's own $v_i$ list --- ``$N$ is the atom's info number.'' If
duplicates are not discarded, $N$ is simply however many entries
remain; the count was never anything other than a property of the
list itself. $A\ \text{Collapse}\ A$ gives $N=2$: two entries, both
reading $A$'s own $(r,\hat n,\chi)$ identically. Shorthand notation
for this case, proposed and worth keeping: $A^{\langle N\rangle}$ for
$N$ identical copies, layered on top of the general bundle notation
rather than replacing it.

\textbf{[UnCollapse's algebra, now fully consistent.]}
\[
A\ \text{UnCollapse}\ A = 0 \quad (N=1\text{ case, matches }\{A\}\ominus A=\emptyset);
\qquad N \to N-1 \text{ per removal, in general.}
\]
$A^{\langle 2\rangle}\ \text{UnCollapse}\ A = A^{\langle 1\rangle} = A$
holds correctly --- specifically because $2-1=1$, not because of any
general collapsing-to-bare-$A$ rule. No remaining contradiction
anywhere in the chain.

\paragraph{What this leaves genuinely open, not resolved here.} The
$\tau$ case --- $A\ \text{Collapse}\ B$ where $A$ and $B$ share
magnitude and position but differ in sort --- is not handled by
$\chi$: $\chi$'s confirmed job is disambiguating multiple things at
one shared direction, nothing about sort, and current Axiom~7
requires every vector in one bundle to share the atom's single
$\tau$. Handling genuinely different-$\tau$ contributions needs the
per-vector $\tau_i$ extension flagged several entries back as not yet
adopted. Not solved by anything in this entry. \emph{[This is now a
formally tracked open item in its own right; see Section~4 above.]}

\section{Considered and found to need nothing filed}

For completeness: the reference-atom and precursor-witness mechanism,
the bundle-vector / uncapped $V_{\text{derived}}$ distinction from
Mix, the $\pi=(x;w_1,\ldots,w_n)$ poem trajectory, and the precise
definition of $\chi$ as internal phase-color (already stated in
Axiom 1: ``two identities may occupy the same direction without
overwriting one another'') were all discussed at length and all
turned out to already be fully present in Inference and Paper 1
respectively. Nothing from those threads is parked here, since there
was nothing left unresolved to park. Checked again against
\texttt{PA\_ReWrite\_V15\_22.tex} specifically: all four are present,
correctly stated, and use current terminology throughout; nothing
here needed a symbol or status correction.

\end{document}
